% !TEX root = chapter-standalone.tex

\section{Upper and lower subsets}
\label{sec:UpperLowerSets}

\todotext{generalize upper/lower set definitions to preorders}

\begin{definition}[Upper set]
    \label{def:upperset}
    An \maindef{upper set}~$\stylesets{U}$ is a subset of a poset~\posA such that, if~$\ela \setin \stylesets{U}$, then all elements of~\posA that are above~$\ela$ are also in~$\stylesets{U}$.
    In other words:
    \begin{equation}
        \prfperiod{\styleelements{x}\setin \stylesets{U}}{\styleelements{x}\posAleq \styleelements{y}}{\styleelements{y}\setin \stylesets{U}}
    \end{equation}
    %In formulas:
    %\begin{equation}
    %  \text{$\stylesets{U}$ is an upperset of $\posA$} \equiv \forall \styleelements{x}\setin \stylesets{S}, \forall \styleelements{y}\setin \posA \colon \styleelements{x}\posAleq \styleelements{y} \Imp \styleelements{y}\setin \stylesets{S}.
    %\end{equation}
\end{definition}
We call~$\setOfUppersets \posA$ the set of \SY{upper sets} of~\posA.

\begin{definition}[Lower set]
    \label{def:lowerset}
    A \maindef{lower set}~$\stylesets{L}$ is a subset of a poset~\posA such that, if~$\ela \setin \stylesets{L}$, then all elements of~\posA that are below~$\ela$ are also in~$\stylesets{L}$.
    In other words:
    \begin{equation}
        \prfperiod{\styleelements{x}\setin \stylesets{L}}{\styleelements{y}\posAleq \styleelements{x}}{\styleelements{y}\setin \stylesets{L}}
    \end{equation}
    %In formulas:
    %\begin{equation}
    %  \text{$S$ is a lower set of $\posA$} \equiv \forall x\setin S, \forall y\setin \posA \colon y\posAleq x \Imp y\setin S.
    %\end{equation}
\end{definition}
We call~$\setOfLowersets\posA$ the set of \SY{lower sets} of~\posA.
%
%\begin{remark}
%  Note that if~$A$ is an antichain of a poset~$\posA$, then the set
%  \begin{equation}
%    I(A)=\makeset{x\colon x\posAleq y, y\setin A}
%  \end{equation}
%  is a lower set of~$\posA$.
%\end{remark}

\begin{marginfigure}
    \subfloat[
        \label{fig:upperset}
    ]{
        \includesag{70_upper_lower_set}
    }
    \\
    \subfloat[
        \label{fig:upperset_b}
    ]{
        \includesag{70_upper_lower_set_b}
    }
    \caption{}
\end{marginfigure}
%

Given the battery choices~$\makeset{\tupp{\unit[10]{\standardcurrency},\unit[500]{g}}, \tupp{\unit[20]{\standardcurrency},\unit[250]{g}}}$, we can represent an \SY{upper set} as in~\cref{fig:upperset}.
The \SY{upper set} can be interpreted as all the potential battery choices which are dominated by at least one of the two choices we have (in case we want to minimize mass and cost).
Similarly, the \SY{lower set} in \cref{fig:upperset_b} can be interpreted as all the potential battery choices which dominate at least one of the choices we have.
Here when considering ``the choices we have'' in \cref{fig:upperset_b}, we not only consider the two choices directly presented to us, but also any convex combination of them.

%\begin{example}[Upper and lower sets in~\Bool]
%  The booleans~$\makeset{\false, \true }$ form a \SY{poset} with~$\false \leq \true\colon(\Bool,\posleq)$ . The subset~$\makeset{\false} \setsubseteq \Bool$ is not an upper set, since~$\false \leq \true$ and~$\true \notsetin \makeset{\false }$.
%\end{example}

\begin{gradedexercise}[\exname{UpperSetsOfPreferences}]
    \label{ex:UpperSetsOfPreferences}

    Consider the poset $\posA$ described by the following Hasse diagram:
    \begin{center}
        \includesag{poset-service}
    \end{center}
    Your task in this exercise is to compute all upper sets of~$\posA$.
\end{gradedexercise}

\solutionof{UpperSetsOfPreferences}

\subsection{Upper and lower closure}

\todotext{decide if DP or CT material; if the latter, generalize to preorders}

\todotext{J: maybe move this to monotone maps section as two nice examples?}


\linkvideo{spring2021-design:up-low-closure} % Upper and lower closure
\begin{definition}[Upper closure operator]
    \label{def:upperclosure}
    The \maindef{upper closure operator}~$\upit $ maps a subset of a poset to the smallest \SY{upper set} that includes it:
    \begin{equation}
        \defmapperiod{
            \upit
        }{
            \posPA
        }{
            \to
        }{
            \uppersets(\posA)
        }{
            \subA
        }{
            \makeset{\styleelements{y}\setin \posAset \mid \exists \styleelements{x}\setin \subA \colon \styleelements{x}\posAleq \styleelements{y}}
        }
    \end{equation}
\end{definition}
\begin{remark}
    Note that, by definition, an \SY{upper set} is closed to \SY{upper closure}.
\end{remark}
\begin{lemma}
    For any~$\setA\setin \posPA$,~$\upit \setA$ is in fact an \SY{upper set}.
\end{lemma}
\begin{proof}
    Suppose~$\styleelements{y}\setin \upit \setA$ and~$\styleelements{z}\setin \posAset$, and suppose $\styleelements{y}\posAleq \styleelements{z}$.
    By definition there exists a~$\styleelements{x}$ such that~$\styleelements{x}\posAleq \styleelements{y}$, meaning that~$\styleelements{x}\posAleq \styleelements{z}$.
    Thus,~$\styleelements{z}\setin \upit \setA$, as was to be shown.
\end{proof}


In the example of battery choices (in the numerical case), first, consider the \SY{upper closure} of a single element of the poset, for instance~$\posAnel{1}=\tupp{\unit[10]{\standardcurrency},\unit[500]{g}}$ (\cref{fig:upperclosure_1}, left).
Second, we can look at the \SY{upper closure} when we add the choice~$\posAnel{2}=\tupp{\unit[20]{\standardcurrency},\unit[250]{g}}$ (\cref{fig:upperclosure_1}, center).

Note that the \SY{upper set} of the subset formed by the two elements is the union of the \SY{upper sets} of the single elements.
%
Finally, we can also define the set
%
\begin{equation}
    \stylesets{S}=\makeset{
        \tupp{\text{cost},\text{mass}}\mid \text{mass}=750-25\cdot \text{cost},
        \forall \text{ cost} \setin [0,20]
    },
\end{equation}
%
and find its \SY{upper closure} (\cref{fig:upperclosure_1}, right).
%
\begin{figure*}[h!]
    \centering
    \includesag{70_uc_1_1}
    \hfill
    \includesag{70_uc_1_2}
    \hfill
    \includesag{70_uc_1_3}
    \caption{Example of \SY{upper closure} for different sets of battery choices.}
    \label{fig:upperclosure_1}
\end{figure*}

\begin{definition}[Lower closure operator]
    \label{def:lowerclosure}
    The \maindef{lower closure operator}~$\downit$ maps a subset to the smallest \SY{lower set} that includes it:
    \begin{equation}
        \defmapperiod{
            \downit
        }{
            \posPA
        }{
            \to
        }{
            \setOfLowersets \posA
        }{
            \stylesets{S}
        }{
            \makeset{ \styleelements{y}\setin \posAset \mid \exists \styleelements{x}\setin \stylesets{S} \colon \styleelements{y}\posAleq \styleelements{x}}
        }
    \end{equation}
\end{definition}



Consider the battery example, and the \SY{antichain} given by the battery models~$\posAnel{1}=\tupp{\unit[10]{\standardcurrency},\unit[1000]{g}}$,~$\posAnel{2}=\tupp{\unit[20]{\standardcurrency},\unit[500]{g}}$, and~$\posAnel{3}=\tupp{\unit[30]{\standardcurrency},\unit[250]{g}}$ (\cref{fig:examplebatt}, left).
The \SY{lower closure} operator~$\downit\makeset{\posAnel{1},\posAnel{2},\posAnel{3}}$ represents all the battery models which, if existing, would dominate~$\makeset{\posAel_1,\posAel_2,\posAel_3}$.
We could instead consider linear maps between the points getting a poset~\posA, and obtain the \SY{lower closure} depicted in \cref{fig:examplebatt} on the right.

\begin{figure*}[h!]
    \centering
    \includesag{70_battery_0}
    \hfill
    \includesag{70_battery_1}
    \hfill
    \includesag{70_battery_1_bis}
    \caption{Example of lower closures.}
    \label{fig:examplebatt}
\end{figure*}

